%% fixed-point-forcing.tex -- two passes of a self-conditioned denoiser, the
%% conventional one and the rollout-derived one, side by side.
%%
%% Panel (b) repeats panel (a) and fades everything it does not change, so the
%% eye lands on the loop and the carry it produces. The fading is a transparency
%% group around the repeated parts, which keeps overlapping strokes from showing
%% through each other.
%%
%% Requires texfig.sty. Input inside a figure* environment:
%%   \input{figures/fixed-point-forcing}
\providecommand{\tffigwidth}{\linewidth}
\resizebox{\tffigwidth}{!}{%
\begin{tikzpicture}[font=]

% ===================================================================
% (a) conventional self-conditioning
% ===================================================================
\node[tflabel, font=\large] (ea) at (0,0) {$\varnothing$};
\node[tfblock={appleBlue}{1.7cm}, minimum height=1.05cm] (da1) at (2.0,0)
      {$D_\theta^t$};
\node[tflabel, font=\large] (sa) at (4.3,0) {$s_t^0$};
\node[tfblock={appleBlue}{1.7cm}, minimum height=1.05cm] (da2) at (7.4,0)
      {$D_\theta^t$};
\node[tflabel, font=\large] (oa) at (9.7,0) {$\hat{x}_1^1$};
\node[tflabel, font=\large] (xa1) at (2.0,2.1) {$x_t$};
\node[tflabel, font=\large] (xa2) at (7.4,2.1) {$x_t$};
\node[tflabelc=appleOrange] (sga) at (5.85,1.0) {Stopgrad};

\draw[tfarrow, draw=appleDark] (ea.east)  -- (da1.west);
\draw[tfarrow, draw=appleDark] (da1.east) -- (sa.west);
\draw[tfarrow, draw=appleDark] (sa.east)  -- (da2.west);
\draw[tfarrow, draw=appleDark] (da2.east) -- (oa.west);
\draw[tfarrow, draw=appleDark] (xa1.south) -- (da1.north);
\draw[tfarrow, draw=appleDark] (xa2.south) -- (da2.north);
\tfcutmark[appleOrange]{(5.85,0)}
\node[tflabel, font=\normalsize] at (4.85,-1.7) {(a) Conventional self-conditioning};

% ===================================================================
% (b) fixed-point forcing: the same pass, rolled out k times
% ===================================================================
\begin{scope}[xshift=13cm]
  \begin{scope}[transparency group, opacity=0.33]
    \node[tflabel, font=\large] (eb) at (0,0) {$\varnothing$};
    \node[tfblock={appleBlue}{1.7cm}, minimum height=1.05cm] (db2) at (7.4,0)
          {$D_\theta^t$};
    \node[tflabel, font=\large] (ob) at (9.7,0) {$\hat{x}_1^{k+1}$};
    \node[tflabel, font=\large] (xb1) at (2.0,2.1) {$x_t$};
    \node[tflabel, font=\large] (xb2) at (7.4,2.1) {$x_t$};
    \node[tflabelc=appleOrange] (sgb) at (5.85,1.0) {Stopgrad};
    \draw[tfarrow] (eb.east) -- (2.0-0.85,0);
    \draw[tfarrow] (4.3+0.45,0) -- (db2.west);
    \draw[tfarrow] (db2.east) -- (ob.west);
    \draw[tfarrow] (xb1.south) -- (2.0,0.53);
    \draw[tfarrow] (xb2.south) -- (db2.north);
    \tfcutmark[appleOrange]{(5.85,0)}
  \end{scope}
  \node[tfblock={appleBlue}{1.7cm}, minimum height=1.05cm] (db1) at (2.0,0)
        {$D_\theta^t$};
  \node[tflabel, font=\large] (sb) at (4.3,0) {$s_t^k$};
  \draw[tfarrow, draw=appleDark] (db1.east) -- (sb.west);
  % the rollout: the carry is fed back k times before the pass that trains
  \draw[tfarrow, draw=appleDark, tfelbow] (sb.north) -- ++(0,1.05)
        -| ($(db1.north east)+(-0.4,0)$);
  \node[tflabel, font=\large\bfseries] at (3.15,2.0) {$\times k$};
  \node[tflabel, font=\normalsize] at (4.85,-1.7) {(b) Fixed-Point Forcing};
\end{scope}

\node[tflabel, font=\itshape] at (11.4,-3.1)
      {Fixed-Point Forcing replaces one-pass carries with rollout-derived carries.};

\end{tikzpicture}}
